% ===================================================================
%  도표 제 9 호 — 산. 숲.
%
%  Traduction, faite en 2026, en coreen d'aujourd'hui, sur l'ido de
%  texto/io. Regles de la colonne : voir 10-dopyo-01.tex. Deux
%  scenes, deux numerotations : les cles portent c1 et c2.
%
%  LE MEME TABLEAU QUE LA COLONNE TELOUGOUE, ET LE RESULTAT INVERSE.
%  Le tableau 9 telougou avait du emprunter neuf noms d'un coup —
%  ఫర్ le sapin, బీచ్ le hetre, ఓక్ le chene, ఎల్మ్ l'orme, పాప్లర్
%  le peuplier, హీథర్ la bruyere, ఫెర్న్ la fougere, ప్రిమ్‌రోజ్ la
%  primevere, మాగ్‌పై la pie — parce qu'aucun de ces vivants ne
%  passe en Andhra. La Coree est a la latitude de la France, et la
%  liste se retourne d'un bloc :
%
%    너도밤나무 (23) le hetre      참나무 (29)   le chene
%    느릅나무 (47) l'orme          자작나무 (26) le bouleau
%    미루나무 (48) le peuplier     전나무 (11)   le sapin
%    고사리 (55)  la fougere       앵초 (9)      la primevere
%    까치 (56)    la pie           노루 (20)     le chevreuil
%    사슴벌레     le lucane        딱따구리 (30) le pic-vert
%    다람쥐 (61)  l'ecureuil       민달팽이 (8)  la limace
%
%  Quatorze noms coreens la ou le telougou en avait neuf d'emprunt.
%  Ce n'est pas une langue plus riche que l'autre : c'est la meme
%  regle — le mot est la ou la chose est — appliquee a un autre
%  parallele. Et le coreen mange la fougere, ce que le livret ne dit
%  pas et que le mot 고사리 dit tout seul.
%
%  QUATRE MANQUENT, ET CE SONT LES QUATRE QU'ON ATTEND : 샤무아 le
%  chamois (54), 히스 la bruyere (57), 마멋 la marmotte (39), 빙카 la
%  pervenche (10). Pyrenees, Alpes, sous-bois d'Europe occidentale.
%  La coupure ne suit pas la langue, elle suit la carte.
%
%  ET LA LIEUE (bloc c1-15-1) TOMBE JUSTE, comme l'ఆమడ telougou. Le
%  coreen mesure les longues marches en 리, et 십 리 fait quatre
%  kilometres — exactement la lieue du livret. Le mot vit encore dans
%  le proverbe : 십 리도 못 가서 발병 난다.
%
%  LE « cigano » (41), QUATRIEME FOIS. Comme au tableau 9 telougou :
%  le livret designe un homme par son peuple, et le mot francais de
%  1926 est aujourd'hui une insulte. La colonne ecrit 떠돌이 사람,
%  « un homme qui va de lieu en lieu ». La scene ne perd rien — ce
%  qui la fait tenir est le contraste du bonnet de fourrure et du
%  capuchon dechire, et il est intact.
%
%  TRENTE-CINQ PAIRES. Le tableau decrit un paysage par attachements
%  — la crete DE la chaine, le nid PRES du fort, la mousse OU le
%  pic-vert cherche, la scie DU bucheron — et le compte s'en ressent.
% ===================================================================

%%K t09-apar-1 apar
\VUsaut{4.0mm}
\VUtitre{15.0pt}{60}{도표 제 9 호}
\VUsaut{3.0mm}

%%K t09-tit-1 sub
\VUcentre{12.0pt}{0}{\textit{첫째 마당.}}
\VUsaut{3.0mm}

%%K t09-tit-2 sub
\VUpk{11.4pt}{40}{산. --- 온천 마을.}
\VUsaut{3.0mm}

%%K t09-c1-01-1 p
1. --- 여드레 전부터 나는 프라츠에 있다. 피레네의 놀라운 \VUgras{산맥} \textsuperscript{(1)}
앞에 섰을 때 느낀 감탄을 다 말할 수가 없다. 그 산맥의 \VUgras{능선} \textsuperscript{(2)}이
푸른 하늘 위에 거대한 꽃장식처럼 도드라진다.

%%K t09-c1-02-1 p
2. --- 피레네의 \VUgras{봉우리} \textsuperscript{(3)}가 그토록 높은 \VUgras{높이}에 이르는
줄은 몰랐다. 훌륭한 \VUgras{빙하} \textsuperscript{(5)} 앞에서, 눈 덮인 \VUgras{뾰족봉}
\textsuperscript{(4)} 앞에서 나는 넋을 잃고 서 있었다. 그 위에서 그토록 무서운
\VUgras{눈사태}가 굴러떨어진다.

%%K t09-tit-3 sub
\VUsaut{3.0mm}
\VUpk{10.2pt}{40}{산 경치.}
\VUsaut{3.0mm}

%%K t09-c1-03-1 p
3. --- 도착한 이튿날 바로, 프라츠 가까이에 있는 꺼진 옛 \VUgras{화산} \textsuperscript{(6)}
에 갔다. 메마르고 헐벗은 \VUgras{분화구} 가장자리에서는, 몇 백 년 전에 \VUgras{산비탈}
\textsuperscript{(7)}로 흘러내린 \VUgras{용암}의 자취가 아직도 또렷이 보인다. 나는 그
용암 조각 몇 개를 \VUgras{비탈} 하나에서 찾아낼 수 있었다.

%%K t09-c1-04-1 p
4. --- 프라츠는 국경에 자리 잡고 있다. 여기 \VUgras{요새} \textsuperscript{(8)}는 —
프랑스와 스페인을 가르는 \VUgras{고갯길} \textsuperscript{(27)}을 지키는 것이 바로 그것이다
— 제 바위 꼭대기에서 \VUgras{먹이}를 노리는 듯한 라인 강가의 옛 성들을 그대로 떠올리게
한다.

%%K t09-c1-05-1 p
5. --- 커다란 \VUgras{독수리} \textsuperscript{(9)} 한 마리가 자주 내 눈길을 끈다. 그
둥지는 \VUgras{요새} \textsuperscript{(8)}에서 멀지 않다.

%%K t09-c1-05-1 p suite
그 \VUgras{맹금}은 \VUgras{부리}가 힘세다. 제 \VUgras{발톱}으로 새끼양이나 새끼염소를
움켜쥐고 제 새끼들에게 자주 물어다 준다. 그 \VUgras{날개}가 어찌나 큰지, 그 무거운 짐을
지고도 오래 하늘에 떠 있을 수 있다.

%%K t09-tit-4 sub
\VUsaut{3.0mm}
\VUpk{10.2pt}{40}{나들이하는 사람.}
\VUsaut{3.0mm}

%%K t09-c1-06-1 p
6. --- 여기에서 가장 즐거운 걸음은 「카우」 \VUgras{폭포} \textsuperscript{(10)}까지 가는
나들이다. \VUgras{떨어지는 물}은 소용돌이치며 떨어졌다가, 고운 \VUgras{시내}가 되어
마을 서쪽 \VUgras{전나무 숲} \textsuperscript{(11)}에서 사라진다. 우리는 그 숲을 지나
\VUgras{기상 관측소} \textsuperscript{(12)}까지 올라갔다. \VUgras{전망대} 꼭대기에서 그
고장의 \VUgras{전경}을 한눈에 보았다. \VUgras{경치}는 놀랍고, \VUgras{자연}은 제가 가진
보물을 이 땅에 아낌없이 쏟아부은 듯하다.

%%K t09-c1-07-1 p
7. --- 내려오려고 우리는 \VUgras{케이블 철도} \textsuperscript{(13)}를 탔고, 옛
\VUgras{성채} \VUgras{옛터} \textsuperscript{(14)} 가까이 작은 \VUgras{정거장}에서 내렸다.
프라츠의 영주들이 예전에 살던 성이다.

%%K t09-c1-08-1 p
8. --- 딱한 성이여! 이제 유행하는 \VUgras{온천장}이 된 저 맵시 있는
\VUgras{온천 마을} \textsuperscript{(15)}을 발밑에 보면서 무슨 생각을 할까?

%%K t09-c1-08-2 p
\VUgras{기후}가 그토록 좋고 \VUgras{광천수}가 그토록 잘 들어서, \VUgras{요양원}
\textsuperscript{(17)}에서 받는 \VUgras{치료}는 참으로 즐거움이다.

%%K t09-tit-5 sub
\VUsaut{3.0mm}
\VUpk{10.2pt}{40}{즐거운 온천 마을.}
\VUsaut{3.0mm}

%%K t09-c1-09-1 p
9. --- 게다가 머무는 것을 즐겁게 하려고 온갖 것을 다 했다. 철길을 놓을 수 있도록 큰
\VUgras{구름다리} \textsuperscript{(16)}를 세웠다. \VUgras{온천장 건물}의
\VUgras{공원} \textsuperscript{(18)}은 — 거기에서 \VUgras{음악회}가 열린다 — 아주
아름답게 꾸며졌다. 고운 「\VUgras{별장}」 \textsuperscript{(19)} 몇 채를 지었는데,
\VUgras{카지노} \textsuperscript{(21)} 둘레에 지었다. 잘 손질된 \VUgras{큰길}
\textsuperscript{(22)}이 오가기를 아주 쉽게 해 준다.

%%K t09-c1-10-1 p
10. --- 수수한 \VUgras{둑} \textsuperscript{(23)} 하나가 \VUgras{길} \textsuperscript{(20)}을
진짜 \VUgras{골짜기} \textsuperscript{(25)}에서 갈라놓는다. 그 골짜기 바닥에는 고운 작은
\VUgras{못} \textsuperscript{(26)}이 있는데, 그것이 맑은 \VUgras{시내} \textsuperscript{(24)}에
물을 댄다.

%%K t09-c1-11-1 p
11. --- 골짜기 건너편에서는 \VUgras{철광} \textsuperscript{(28)}이 돌아간다. 나는
\VUgras{광부}들이 \VUgras{수직갱}으로 내려가는 것을 보러 여러 번 갔고, 그들이 하루를
보내는 \VUgras{갱도}에도 들어가 보았다. 거기에서 그들은 \VUgras{쇠붙이}가 든
\VUgras{광석}을 캐낸다.

%%K t09-c1-12-1 p
12. --- 광산 가까이에 큰 \VUgras{제련소}를 지었다. 거기에서 캐낸 값진 것을 곧바로
쓰려는 것이다. 여기에 넉넉한 \VUgras{석탄}으로 데우는 \VUgras{용광로} \textsuperscript{(29)}
덕분에 \VUgras{무쇠}를 빠르고 싸게 얻는다.

%%K t09-c1-13-1 p
13. --- 광산에서 멀지 않은 곳에 \VUgras{채석장} \textsuperscript{(30)}을 판다. 늙은
\VUgras{석수}가 제 \VUgras{곡괭이}로 떼어 내는 것은 \VUgras{화강암}도 \VUgras{반암}도
\VUgras{사암}도 아니고, 집 짓는 데 쓰는 수수한 \VUgras{다듬돌}이다.

%%K t09-c1-14-1 p
14. --- 이 고장의 아름다운 꾸밈 가운데 하나는 \VUgras{전나무} \textsuperscript{(31)}다.
어디에서나 — \VUgras{후미진 골} \textsuperscript{(32)} 바닥에, \VUgras{급류}
\textsuperscript{(33)} 기슭에, \VUgras{심연} \textsuperscript{(34)} 가장자리에, 또는
낭떠러지 곁에 — 그 가느다란 \VUgras{바늘잎}이 푸른 빛깔을 놓는다.

%%K t09-tit-6 sub
\VUsaut{3.0mm}
\VUpk{10.2pt}{40}{걷기.}
\VUsaut{3.0mm}

%%K t09-c1-15-1 p
15. --- 나는 방학을 \VUgras{멀리 걷는} 데 쓴다. 산을 굽이굽이 도는 \VUgras{길}
\textsuperscript{(35)}은 거리를 느끼지 못한 채 여러 \VUgras{킬로미터}, 자주 여러
\VUgras{십 리}를 걷게 해 준다.

%%K t09-c1-15-2 p
\VUgras{이정표} \textsuperscript{(36)}와 \VUgras{길 알림 기둥} \textsuperscript{(37)}이 그
길에 표를 놓는다. 그 길에서는 젊은 \VUgras{산골 사람} \textsuperscript{(38)}을 자주
만난다 — 옆구리에 \VUgras{걸망} \textsuperscript{(40)}을 차고, \VUgras{마멋}
\textsuperscript{(39)}을 안고 간다. 또는 어떤 \VUgras{떠돌이 사람} \textsuperscript{(41)}을
만나는데, 그의 \VUgras{털모자} \textsuperscript{(42)}가 찢어진 낡은 \VUgras{두건}
\textsuperscript{(43)}과 묘하게 어긋난다. 그는 사슬로 길들인 \VUgras{곰} \textsuperscript{(44)}
을 끌고 간다.

%%K t09-tit-7 sub
\VUpk{10.2pt}{40}{산 오르기.}
\VUsaut{3.0mm}

%%K t09-c1-16-1 p
16. --- 거의 날마다 나는 \VUgras{길잡이} \textsuperscript{(48)}와 함께 \VUgras{산 오르기}를
익힌다. 등에 \VUgras{배낭} \textsuperscript{(47)}을 지고 손에 「\VUgras{지팡이}」를 들고
참된 \VUgras{나그네} \textsuperscript{(46)}처럼 \VUgras{바위} \textsuperscript{(45)}에 달려들
때 나는 아주 자랑스럽다.

%%K t09-c1-17-1 p
17. --- 산꼭대기에서 보면 벌판 사람들이 개미보다 크게 보이지 않는다. \VUgras{양떼}
\textsuperscript{(50)}는 — \VUgras{목장} \textsuperscript{(49)}에서 풀을 뜯는다 — 아이들의
장난감처럼 보인다. \VUgras{소나무} \textsuperscript{(51)}나 \VUgras{통나무집}
\textsuperscript{(52)}도 그 푸르름 위의 한 점 얼룩으로밖에 보이지 않는다.

%%K t09-c1-18-1 p
18. --- 그 나들이에서 우리는 \VUgras{오솔길} \textsuperscript{(53)} 위에서
\VUgras{샤무아 사냥꾼} \textsuperscript{(54)}을 자주 만난다 — 방금 잡은 짐승을 어깨에
메고 가는 사람이다.

%%K t09-c1-19-1 p
19. --- 지난번 걸음에서 나는 아주 볼 만한 \VUgras{고사리} \textsuperscript{(55)} 한 잎과
고운 분홍빛 \VUgras{히스} \textsuperscript{(57)} 한 가지를 내 식물첩에 넣었다. 그것들은 작은
\VUgras{굴} \textsuperscript{(56)} 어귀에서 딴 것이다.

%%K t09-tit-8 sub
\VUsaut{3.0mm}
\VUcentre{12.0pt}{0}{\textit{둘째 마당.}}
\VUsaut{3.0mm}

%%K t09-tit-9 sub
\VUpk{11.4pt}{40}{숲. --- 사냥.}
\VUsaut{3.0mm}

%%K t09-c2-01-1 p
1. --- 어제 나는 숲에서 아주 즐거운 하루를 보냈다. 거기에 사는 짐승들과 \VUgras{벌레}
\textsuperscript{(5)}들 사이의 부지런함이 나에게 큰 흥미를 주었다.

%%K t09-c2-01-2 p
\VUgras{거미} \textsuperscript{(1)}는 — 제 \VUgras{거미줄} \textsuperscript{(2)}을 두 가지
사이에 친다, 지나가는 \VUgras{파리} \textsuperscript{(3)}를 잡으려고 —, 제 껍데기를 힘겹게
지고 가는 \VUgras{달팽이} \textsuperscript{(4)}, 제 먹이를 살피는 \VUgras{사슴벌레},
\VUgras{개미집} \textsuperscript{(6)} 가까이에서 그토록 열심인 부지런한 \VUgras{개미} —
그 작은 것들이 모두 남다른 부지런함으로 가고 오고 일했다.

%%K t09-c2-02-1 p
2. --- \VUgras{뱀} \textsuperscript{(7)} 한 마리가 나무 밑동 아래 몸을 사리고 있었다.
처음 보고는 \VUgras{살무사}인 줄 알았으나, 수수한 \VUgras{무자치}에 지나지 않았다.
때때로 가느다란 머리를 내밀었는데, 그 머리는 언제나 물 채비가 된 듯이 보였다. 내 앞에서
굵은 \VUgras{민달팽이} \textsuperscript{(8)} 한 마리가 무겁게 기어갔다 — 거기에 그득히
자라는 맛있는 \VUgras{산딸기} \textsuperscript{(11)} 하나를 먹어 치운 뒤에.

%%K t09-c2-03-1 p
3. --- 땅바닥은 \VUgras{숲꽃}으로 깔려 있었다. \VUgras{앵초} \textsuperscript{(9)},
\VUgras{빙카} \textsuperscript{(10)}, 그리고 내가 모르는 많은 풀이 \VUgras{풀숲}
\textsuperscript{(12)} 사이에 피어 있었다.

%%K t09-c2-04-1 p
4. --- 그 고장에서는 송로버섯이 \VUgras{버섯} \textsuperscript{(14)}만큼이나 흔하다.
산지기의 \VUgras{새끼 돼지} \textsuperscript{(13)} 한 마리가, 그것을 몇 개 찾아낼 수 있을까
하고 게걸스레 뒤졌다.

%%K t09-tit-10 sub
\VUsaut{3.0mm}
\VUpk{10.2pt}{40}{밀렵.}
\VUsaut{3.0mm}

%%K t09-c2-05-1 p
5. --- 나는 아주 재미있는 한 장면의 \VUgras{끝}을 보았다. 제 \VUgras{바셋 사냥개}
\textsuperscript{(15)}의 코 덕분에 \VUgras{산지기} \textsuperscript{(16)}가 방금
\VUgras{밀렵꾼} \textsuperscript{(22)}을 붙잡았다 — 그가 잡은 \VUgras{사냥감}
\textsuperscript{(17)}을 메고 가려던 바로 그 순간에. 붙잡히느니 제 노획물을 버리는 편이
낫다고 여겨 밀렵꾼은 달아났고, \VUgras{산토끼} \textsuperscript{(18)},
\VUgras{암사슴} \textsuperscript{(19)}, \VUgras{노루} \textsuperscript{(20)},
\VUgras{멧돼지} \textsuperscript{(21)}가 풀 위에 놓여 산지기를 기쁘게 했다.

%%K t09-c2-06-1 p
6. --- 숲이 품은 나무의 가짓수는 놀랍다. \VUgras{너도밤나무} \textsuperscript{(23)}와
\VUgras{자작나무} \textsuperscript{(26)}, \VUgras{소나무} — 이것은 \VUgras{송진}
\textsuperscript{(27)}을 낸다 —, \VUgras{참나무} \textsuperscript{(29)}, 그리고 다른 많은
나무가 제 가지를 뒤섞는다.

%%K t09-c2-06-2 p
\VUgras{담쟁이덩굴} \textsuperscript{(24)}은 키 큰 나무의 밑동에 붙어, 반짝이는 푸르름으로
\VUgras{큰 나무 숲} \textsuperscript{(25)}에 빛깔을 준다. 그 푸르름은 \VUgras{이끼}
\textsuperscript{(31)}의 더 옅고 창백한 빛깔과 기분 좋게 어울린다 — 그 이끼 속에서
\VUgras{딱따구리} \textsuperscript{(30)}가 벌레를 찾는다.

%%K t09-tit-11 sub
\VUsaut{3.0mm}
\VUpk{10.2pt}{40}{숲 경치.}
\VUsaut{3.0mm}

%%K t09-c2-07-1 p
7. --- 늙은 참나무들이 나를 사로잡았다. 땅에서 반쯤 나온 그 굵고 비틀린 \VUgras{뿌리}
\textsuperscript{(32)}가 나무에 힘과 기운의 훌륭한 모습을 준다. 그래서 나는
\VUgras{땅임자} \textsuperscript{(35)}가 어떻게 그것을 베게 하고 \VUgras{톱}
\textsuperscript{(34)}에 넘길 수 있는지 스스로 묻는다 — \VUgras{나무꾼} \textsuperscript{(33)}
의 톱이다. 딱한 \VUgras{밑동} \textsuperscript{(36)}들이, \VUgras{껍질} \textsuperscript{(37)}
이 벗겨지거나 \VUgras{도끼} \textsuperscript{(38)}에 쪼개진 밑동들이 — \VUgras{날품팔이}
\textsuperscript{(39)}의 도끼다 — 나를 안타깝게 한다.

%%K t09-c2-08-1 p
8. --- 그 곁에서 늙은 \VUgras{숯 굽는 이} \textsuperscript{(41)}가 제 \VUgras{삽}으로 숯
\VUgras{더미} \textsuperscript{(42)}를 채비하고 있었고, 한 노파가 베어 낸 나무의 잔가지로
\VUgras{마른 나무} \VUgras{단} \textsuperscript{(40)}을 지고 갔다.

%%K t09-tit-12 sub
\VUsaut{3.0mm}
\VUpk{10.2pt}{40}{사냥.}
\VUsaut{3.0mm}

%%K t09-c2-09-1 p
9. --- 잠시 쉬려고 \VUgras{산지기 집} \textsuperscript{(46)}에 가려 했다. 그러나
\VUgras{작은 못} \textsuperscript{(43)} 가까이에 이르렀을 때 — 그 위로 고운 \VUgras{백조}
\textsuperscript{(45)}가 헤엄친다 — 요즘 난 \VUgras{물난리} \textsuperscript{(44)}가 길을
진짜 \VUgras{늪}으로 바꾸어 놓은 것을 알아차렸다. 나는 돌아가야 했다. 숲 가장자리에 선
\VUgras{개잎갈나무} \textsuperscript{(49)}와 \VUgras{미루나무} \textsuperscript{(48)}와
\VUgras{느릅나무} \textsuperscript{(47)} 곁을 지나면서, \VUgras{덤불} \textsuperscript{(54)}
에서 훌륭한 \VUgras{수사슴} \textsuperscript{(50)} 한 마리가 뛰쳐나오는 것을 보았다 —
끈질긴 \VUgras{사냥개 떼} \textsuperscript{(51)}가 그것을 쫓고 있었고, 그 개들을
\VUgras{뿔피리} \textsuperscript{(53)}도 부추기고 있었다 — \VUgras{말 탄 사냥꾼}
\textsuperscript{(52)}들의 뿔피리다. 딱한 짐승은 아주 지쳐 있었다. 그것은 내게서 몇
걸음 떨어진 곳에 죽어 가며 쓰러졌다.

%%K t09-c2-10-1 p
10. --- \VUgras{몰이사냥}의 소란에 끌려 산지기의 아들이 집에서 나왔고, 우리 둘은 다시
숲으로 들어갔다. 그는 늙은 참나무 가지 위에서 \VUgras{까치} \textsuperscript{(56)}를
보았다. 그 둥지가 \VUgras{잎사귀} \textsuperscript{(58)} 속에 숨어 있으리라 여기고, 그
둥지를 잡으려 했다.

%%K t09-c2-10-2 p
\VUgras{다람쥐} \textsuperscript{(61)}처럼 날래게 그는 \VUgras{가지} \textsuperscript{(55)}
에서 가지로 올라갔다. 그러나 아무것도 찾지 못하고, 늙은 \VUgras{까마귀} \textsuperscript{(60)}
한 마리를 날아오르게 한 것이 고작이었다 — 그 까마귀는 \VUgras{잔가지} \textsuperscript{(57)}
에 앉아 있었다.

\VUsaut{6.0mm}
\VUfilet{22mm}
